\input{../tex-inputs/latex-leseansicht-vorspann}

\section[Arthur Schnitzler an Gustav Schwarzkopf, 23. 6. 1899]{L04437 Arthur Schnitzler an Gustav Schwarzkopf, 23. 6. 1899}
\nopagebreak\mylabel{L04437v}
\rehead{ }\normalsize\beginnumbering\briefempfaengerindex{Schwarzkopf, Gustav@\textsc{Schwarzkopf, Gustav}!zzzSchnitzler, Arthur@\emph{von Arthur Schnitzler}!1899-06-231@{23. 6. 1899}|(be}
\toendnotes[C]{\smallbreak\pagebreak[2]}
\correspDesc{Versand  durch Arthur Schnitzler am 23. 6. 1899 in Belišće
\newline{}Übermittlung  am 23. 6. 1899 in Belišće
\newline{}Erhalt  durch Gustav Schwarzkopf im Zeitraum [24. 6. 1899 – 28. 6. 1899?] in Tiefer Graben 23}\toendnotes[C]{\smallbreak}
\Standort{DLA, A:Schnitzler, HS.1985.1.1897.}
\physDesc{Bildpostkarte, 101 Zeichen
\newline{}Handschrift: Bleistift, deutsche Kurrent
\newline{}Versand: Stempel: »\nobreak{}\oindex{Belišće@\textbf{Belišće}|pwk}Belišće, 99 Jun {[}2{]}3\nobreak{}«.  
\newline{}Zusatz: mit schwarzer Tinte von unbekannter Hand Ergänzung der Diakritika beim gedruckten
                                 Ortsnamen: »Belišće\oindex{Belišće@\textbf{Belišće}|pw}« }\pstart{}{\pb}Hrn \textsc{Gustav Schwarzkopf}\pend{}\pstart{}Wien\oindex{Wien@\textbf{Wien}|pw}\pend{}\pstart{}\textsc{I. Tiefer Graben 23}\oindex{Wien@\textbf{Wien}!I., Innere Stadt@\textbf{I., Innere Stadt}!Tiefer Graben 23@\textbf{Tiefer Graben 23}|pw}.\pend{}{\bigskip}
\pstart
           \centering{}{\pb}\textcolor{gray}{\textbf{Sägewerk S. H. Gutmann\oindex{Sägewerk S. H. Gutmann@\textbf{Sägewerk S. H. Gutmann}|pw}.}}\pend
           
\pstart
           \centering{}\textcolor{gray}{\textbf{Gruss aus Belisce\oindex{Belišće@\textbf{Belišće}|pw}.}}\pend
           \vspace{1em}
\pstart
           \noindent{}{\pb}bei ſtrömendem Regen, und herzlich.\pend
           \pstart Ihr \spacefill\mbox{Arth Sch.}\pend{}
\pstart
           23. 6. 99\pend
           \selectlanguage{ngerman}\endnumbering\briefempfaengerindex{Schwarzkopf, Gustav@\textsc{Schwarzkopf, Gustav}!zzzSchnitzler, Arthur@\emph{von Arthur Schnitzler}!1899-06-231@{23. 6. 1899}|)be}\mylabel{L04437h}
\begin{anhang}
\end{anhang}\newcommand{\dateiname}{L04437}\newcommand{\titel}{Arthur Schnitzler an Gustav Schwarzkopf, 23. 6. 1899}\newcommand{\editorInnen}{Herausgegeben von Jahnke, SelmaMüller, Martin Anton}\input{../tex-inputs/latex-leseansicht-abspann}
